\documentclass[10pt,a4paper]{article}

\usepackage[a4paper,margin=13mm]{geometry}
\usepackage{amsmath}
\usepackage{booktabs}
\usepackage{enumitem}
\usepackage{microtype}
\usepackage[table]{xcolor}
\usepackage{tabularx}
\usepackage{titlesec}
\usepackage[hidelinks]{hyperref}

\definecolor{accent}{HTML}{28536B}
\definecolor{light}{HTML}{EAF1F4}
\setlength{\parindent}{0pt}
\setlength{\parskip}{2pt}
\setlist[itemize]{leftmargin=4.5mm,itemsep=0pt,topsep=1pt}
\titleformat{\section}{\color{accent}\bfseries\large}{}{0pt}{}
\titlespacing*{\section}{0pt}{5pt}{2pt}
\pagestyle{empty}

\newcommand{\commit}[1]{\texttt{#1}}
\newcommand{\pathcode}[1]{\texttt{\small #1}}

\begin{document}

\small

{\color{accent}\LARGE\bfseries Linear Asset-Market Teaching Model}\par
{\large Handout on commits \commit{a3a3020} and \commit{1e555f8}}\hfill
{\small 7 August 2026}

\vspace{1mm}
\colorbox{light}{%
  \parbox{\dimexpr\textwidth-2\fboxsep\relax}{%
    \textbf{Outcome.} The teaching note and its canonical figures now use one
    transparent linear asset-market architecture. The former nonlinear
    implementation and figures remain available in archive folders. A follow-up
    change centers the asset-price effect on the reference value $AP=1$.
  }%
}

\section*{1. From the nonlinear to the linear model \hfill\normalfont\commit{a3a3020}}

\begin{itemize}
  \item The linear model became the canonical implementation in
    \pathcode{src/static/models/asset\_model.jl}; the nonlinear version moved to
    \pathcode{src/static/models/archive/asset\_model\_clarifications\_nonlinearapproach.jl}.
  \item The teaching note's equations, descriptions, comparative statics, model
    variants, and captions were aligned with the linear implementation.
  \item The five linear figures were promoted to their canonical filenames.
    Their nonlinear predecessors moved to \pathcode{plots/archive/}.
  \item Plot generation and tests now refer only to the canonical model names:
    \texttt{Baseline}, \texttt{PQC}, \texttt{PQCr}, \texttt{PQCrDIFF}, and
    \texttt{FirmsRation}.
\end{itemize}

The common asset-market block is
\begin{align*}
  SD &= s_0-s_1r-s_2(AP-1),
  & AD &= \gamma_0+\frac{SD}{1-\gamma},\\
  AP &= p_1\frac{AD}{AS},
  & AS &= AQ(\alpha+g_a).
\end{align*}
Here $SD$ is nominal speculative debt, $AD$ is gross nominal asset demand,
$AS$ is the quantity offered for trade, and $AP$ is the asset-price index.
Seller reinvestment produces the within-period multiplier $1/(1-\gamma)$.

\section*{2. Centering the price response at $AP=1$ \hfill\normalfont\commit{1e555f8}}

The initial linear version used $-s_2AP$. The final specification uses
$-s_2(AP-1)$, so $s_0-s_1r$ denotes speculative debt at the reference asset
price. The slope is unchanged: higher prices still reduce desired speculative
borrowing, providing stabilizing rather than self-reinforcing feedback.

Substitution gives the closed-form equilibrium
\[
  AP=\frac{p_1\left[\gamma_0+(s_0+s_2-s_1r)/(1-\gamma)\right]}
           {AS+p_1s_2/(1-\gamma)}.
\]
Centering therefore adds $s_2$ to the numerator while leaving the stabilizing
$s_2$ term in the denominator. The Julia API retains the existing parameter
names \texttt{s1} and \texttt{s2}; economically these are the rate and
asset-price sensitivities $s_r$ and $s_q$.

\section*{3. What distinguishes the five variants}

\renewcommand{\arraystretch}{1.05}
\rowcolors{2}{light}{white}
\begin{tabularx}{\textwidth}{@{}lX@{}}
  \toprule
  \textbf{Variant} & \textbf{Transmission mechanism} \\
  \midrule
  Baseline & The asset market expands loans, deposits, and bank balance sheets,
  but does not affect real activity. \\
  PQC & Gross asset demand reduces the share of productive credit that banks
  accommodate. \\
  PQCr & Adds an asset-price response to the policy rule on top of the PQC
  credit-rationing channel. \\
  PQCrDIFF & Applies a penalty multiplier to the speculative borrowing rate and
  makes credit rationing depend on $SD$. \\
  FirmsRation & Firms reduce productive credit demand directly when asset prices
  rise; bank credit rationing remains exogenous. \\
  \bottomrule
\end{tabularx}

\section*{4. Reproducibility and review}

\begin{itemize}
  \item Independent review confirmed the model equations and every derived
    \texttt{IS}, \texttt{IR}, \texttt{ADc}, and \texttt{AMD} curve formula.
  \item The complete Julia test suite passed: \textbf{335/335 tests}; package
    loading and whitespace checks also passed.
  \item All five canonical asset-model plots were regenerated and visually
    reviewed. The unrelated SimplePK plot remained byte-identical.
  \item The teaching note rebuilt successfully as a 21-page PDF. Remaining
    citation, layout, and duplicate-destination warnings predate these changes.
\end{itemize}

\vfill
{\footnotesize\color{accent}
Repository state after the two commits: the linear specification is canonical,
the nonlinear material is preserved for comparison, and no changes were pushed
to the remote as part of this workflow.}

\end{document}
